% !TEX program = lualatex

% 载入 SJTUThesis 模版
\documentclass[type=master,math-style=TeX]{sjtuthesis}
% 选项
%   type=[doctor|master|bachelor],     % 可选（默认：master），论文类型
%   zihao=[-4|5],                      % 可选（默认：-4），正文字号大小
%   lang=[zh|en|de|ja],                % 可选（默认：zh），论文的主要语言
%   review,                            % 可选（默认：关闭），盲审模式
%   [twoside|oneside],                 % 可选（默认：twoside），双页或单页边距模式
%   [openright|openany],               % 可选（默认：openright），奇数页或任意页开始新章
%   math-style=[ISO|TeX],              % 可选 (默认：ISO)，数学符号样式

\usepackage{sjtutex-markdown}

% 基础设置
\sjtusetupyaml{contents/thesis-setup.yaml}

% 将内容路径加入图像搜索路径，以便插入图片和预览文档
\graphicspath{{contents/}}

% 参考文献设置
\usepackage[backend=biber,style=gb7714-2015]{biblatex}
\addbibresource{ref.bib}

\begin{document}

%TC:ignore

% 标题页
\maketitle

% 原创性声明及使用授权书
\copyrightpage
% 插入外置原创性声明及使用授权书
% 此时必须在导言区使用 \usepackage{pdfpages}
% \copyrightpage[file=scans/sample-copyright-a.pdf]

% 前置部分
\frontmatter

% 摘要
\begin{abstract}[zh]
    \markdownInput{contents/thesis-abstract-zh.md}
\end{abstract}
\begin{abstract}[en]
    \markdownInput{contents/thesis-abstract-en.md}
\end{abstract}

% 目录
\tableofcontents
% 插图索引
\listoffigures
% 表格索引
\listoftables

%TC:endignore

% 主体部分
\mainmatter

% 正文内容
\markdownInput{contents/thesis-mainmatter.md}

%TC:ignore

% 参考文献
\printbibliography[heading=bibintoc]

% 附录
\appendix

% 附录中图表不加入索引
\captionsetup{list=no}

% 附录内容

% 符号对照表
\begin{nomenclature}
  \begin{center}
    \markdownInput{contents/thesis-nomenclature.md}
  \end{center}
\end{nomenclature}

% 结尾部分
\backmatter

% 用于盲审的论文需隐去致谢、发表论文、科研成果、简历

% 致谢
\begin{acknowledgements}
  \markdownInput{contents/thesis-acknowledgements.md}
\end{acknowledgements}

% 发表论文及科研成果
% 盲审论文中，发表论文及科研成果等仅以第几作者注明即可，不要出现作者或他人姓名
\begin{achievements}
  \subsection*{学术论文}
  \begin{bibliolist}{00}
    \item Chen H, Chan C~T. Acoustic cloaking in three dimensions using acoustic metamaterials[J]. Applied Physics Letters, 2007, 91:183518.
    \item Chen H, Wu B~I, Zhang B, et al. Electromagnetic Wave Interactions with a Metamaterial Cloak[J]. Physical Review Letters, 2007, 99(6):63903.
  \end{bibliolist}
  \begin{bibliolist*}{00}
    \item 第一作者. 中文核心期刊论文, 2007.
    \item 第一作者. EI 国际会议论文, 2006.
  \end{bibliolist*}

  \subsection*{专利}
  \begin{bibliolist}{00}
    \item 第一发明人，“永动机”，专利申请号202510149890.0
  \end{bibliolist}
  \begin{bibliolist*}{00}
    \item 第一发明人，“永动机”，专利申请号XXXXXXXXXXXX.X
  \end{bibliolist*}
\end{achievements}

% 学士学位论文要求在最后有一个大摘要，单独编页码
% 如果撰写学士学位论文，请取消下一行的注释
\begin{digest}
  \markdownInput{contents/thesis-digest.md}
\end{digest}

%TC:endignore

\end{document}
